\documentclass[11pt,a4paper]{article}

% --- Encoding and Japanese ---
\usepackage[utf8]{inputenc}
\usepackage[T1]{fontenc}
\usepackage{CJKutf8}

% --- Mathematics ---
\usepackage{amsmath,amssymb,amsthm}

% --- Layout ---
\usepackage[margin=25mm]{geometry}
\usepackage{booktabs}
\usepackage{graphicx}
\usepackage{hyperref}
\usepackage{algorithm}
\usepackage{algpseudocode}
\usepackage{multirow}
\usepackage{url}

% --- Theorem environments ---
\newtheorem{theorem}{定理}
\newtheorem{definition}[theorem]{定義}
\newtheorem{lemma}[theorem]{補題}
\newtheorem{proposition}[theorem]{命題}
\newtheorem{corollary}[theorem]{系}

\begin{document}
\begin{CJK}{UTF8}{ipxm}

% =====================================================================
\title{シングルセルトランスクリプトミクスにおける\\擬時間上の密集遺伝子共発現区間の検出}
\author{匿名}
\date{}
\maketitle

% =====================================================================
\begin{abstract}
シングルセル RNA シーケンシング（scRNA-seq）は細胞レベルの遺伝子発現プロファイリングを可能にするが、
擬時間に沿った遺伝子共発現パターンの時間的局在を検出する手法は未確立である。
本研究では、密集区間マイニング（Dense Interval Mining）フレームワークを scRNA-seq データに適用し、
擬時間上で特定の遺伝子セットが密集的に共発現する区間（Dense Gene Co-expression Interval; DGCI）を
検出する手法を提案する。
発現閾値による二値化とApriori+スライディングウィンドウアルゴリズムを組み合わせ、
既知の分化段階に対応する遺伝子共発現モジュールの時間的局在を自動的に同定する。
造血分化を模した合成 scRNA-seq データを用いた評価では、
5 つの既知モジュール全てを検出（モジュール再現率 1.00）し、
平均 IoU 0.94 の高い時間的精度を達成した。
本手法は、従来の静的共発現解析（WGCNA 等）では捉えられない
時間的に局在した遺伝子協調パターンの発見を可能にする。
\end{abstract}

\noindent\textbf{キーワード:}
シングルセル RNA-seq, 擬時間, 密集区間マイニング, 遺伝子共発現, Apriori アルゴリズム

% =====================================================================
\section{はじめに}
\label{sec:intro}

シングルセル RNA シーケンシング（scRNA-seq）技術の発展により、
個々の細胞レベルでの遺伝子発現プロファイリングが可能となった~\cite{trapnell2014dynamics}。
擬時間（pseudotime）解析は、細胞を連続的な分化軌跡に沿って並べることで、
分化・細胞周期・刺激応答といった生物学的プロセスの動態を捉える手法である~\cite{street2018slingshot,haghverdi2016diffusion}。

しかし、現行の軌跡解析手法には重要な限界がある。
tradeSeq~\cite{vandenberge2020tradeseq}や GeneSwitches~\cite{cao2020geneswitches}は
個々の遺伝子レベルの発現変動を検出するが、遺伝子セットの共発現パターンは扱わない。
一方、WGCNA~\cite{langfelder2008wgcna}や SCENIC~\cite{aibar2017scenic}は
遺伝子共発現ネットワークを構築するが、時間的局在の情報を持たない。
すなわち、「どの遺伝子セットが、擬時間のどの区間で密集的に共発現するか」を
直接検出する手法は存在しない。

本研究では、頻出アイテムセットマイニングのスライディングウィンドウ拡張である
密集区間マイニング（Dense Interval Mining）フレームワークを scRNA-seq データに適用し、
この問題を解決する。具体的な貢献は以下の通りである。

\begin{enumerate}
  \item \textbf{問題定式化}: 擬時間上の密集遺伝子共発現区間（DGCI）の概念を定義する。
  \item \textbf{アダプタ設計}: scRNA-seq 発現行列を密集区間マイニングの入力形式に変換する
        アダプタを設計する。
  \item \textbf{実験的検証}: 造血分化を模した合成データで、既知の遺伝子モジュールの
        検出精度を定量的に評価する。
\end{enumerate}

% =====================================================================
\section{関連研究}
\label{sec:related}

\paragraph{擬時間推定}
Monocle~\cite{trapnell2014dynamics}は逆グラフ埋め込みにより細胞の分化軌跡を推定する。
Slingshot~\cite{street2018slingshot}はクラスタベースの最小全域木上で主曲線を推定する。
DPT~\cite{haghverdi2016diffusion}は k-NN グラフ上のランダムウォーク距離を用いる。
これらは擬時間推定には優れるが、下流の共発現解析機能は限定的である。

\paragraph{遺伝子共発現解析}
WGCNA~\cite{langfelder2008wgcna}はサンプル間の相関を用いて共発現モジュールを同定するが、
時間的局在の概念を持たない。
SCENIC~\cite{aibar2017scenic}は共発現とモチーフ解析を組み合わせて制御ネットワークを推定するが、
擬時間に沿った動的変化は捉えない。

\paragraph{軌跡上の発現解析}
tradeSeq~\cite{vandenberge2020tradeseq}は GAM ベースの軌跡上差異発現解析を提供するが、
個々の遺伝子レベルである。
CellRank~\cite{lange2022cellrank}は RNA velocity と擬時間を統合するが、
共発現パターンの検出は対象外である。

\paragraph{バイオインフォマティクスにおけるアイテムセットマイニング}
Creighton \& Hanash~\cite{creighton2003mining}はマイクロアレイデータに
アソシエーションルールマイニングを適用したが、静的解析であり時間的構造を扱わない。
擬時間に沿った scRNA-seq データへの密集アイテムセットマイニングの適用は
本研究が初めてである。

% =====================================================================
\section{手法}
\label{sec:method}

\subsection{問題定義}

\begin{definition}[擬時間順序付き細胞列]
$\mathcal{C} = \{c_1, c_2, \ldots, c_n\}$ を $n$ 個の細胞とし、
擬時間 $\tau(c_1) \leq \tau(c_2) \leq \cdots \leq \tau(c_n)$ で順序付けられているものとする。
\end{definition}

\begin{definition}[発現二値化]
遺伝子全体 $\mathcal{G} = \{g_1, \ldots, g_m\}$ に対し、
細胞 $c_i$ における遺伝子 $g_j$ の正規化発現値を $x_{ij}$ とする。
閾値 $\theta_j$ を用いた二値化指標を次のように定義する：
\begin{equation}
  b_{ij} = \begin{cases} 1 & x_{ij} \geq \theta_j \\ 0 & \text{otherwise} \end{cases}
\end{equation}
\end{definition}

\begin{definition}[トランザクション]
細胞 $c_i$ のトランザクションを $T_i = \{g_j \in \mathcal{G} \mid b_{ij} = 1\}$ とする。
擬時間順序付きトランザクションデータベースは $\mathcal{D} = (T_1, T_2, \ldots, T_n)$ である。
\end{definition}

\begin{definition}[ウィンドウサポート]
\label{def:window-support}
遺伝子セット $S \subseteq \mathcal{G}$、ウィンドウサイズ $W$、ウィンドウ左端 $l$ に対し：
\begin{equation}
  \mathrm{sup}(S, l, W) = |\{i \mid l \leq i \leq l+W,\; S \subseteq T_i\}|
\end{equation}
\end{definition}

\begin{definition}[密集遺伝子共発現区間 (DGCI)]
\label{def:dgci}
遺伝子セット $S$ とサポート閾値 $\sigma$ に対し、
連続する擬時間インデックス区間 $[s, e]$ が $S$ の密集区間であるとは：
\begin{equation}
  \forall l \in [s, e]:\; \mathrm{sup}(S, l, W) \geq \sigma
\end{equation}
が成立し、かつ $[s,e]$ が極大であることをいう。
組 $(S, [s, e])$ を密集遺伝子共発現区間（DGCI）と呼ぶ。
\end{definition}

\subsection{アルゴリズム}

\begin{theorem}[Apriori 単調性の保存]
\label{thm:apriori-monotonicity}
遺伝子セット $S$ が密集区間を持たないならば、
任意の上位集合 $S' \supset S$ も密集区間を持たない。
\end{theorem}

\begin{proof}
任意のトランザクション $T_i$ に対し、$S \subseteq T_i$ は $S' \subseteq T_i$ の
必要条件である。
従って $\mathrm{sup}(S', l, W) \leq \mathrm{sup}(S, l, W)$ が全ての $l$ に対し成立し、
$S$ が密集条件を満たす区間が存在しなければ $S'$ も満たし得ない。
\end{proof}

この性質により、Apriori アルゴリズムの候補剪定が有効に機能する。
具体的な手順を Algorithm~\ref{alg:dgci} に示す。

\begin{algorithm}[t]
\caption{DGCI マイニング}
\label{alg:dgci}
\begin{algorithmic}[1]
\Require 発現行列 $X$, 擬時間 $\boldsymbol{\tau}$, ウィンドウサイズ $W$, サポート閾値 $\sigma$, 最大サイズ $k_{\max}$
\Ensure DGCI 集合 $\mathcal{R}$
\State 細胞を擬時間でソートし、発現を二値化してトランザクション $\mathcal{D}$ を構築
\State $\mathcal{F}_1 \gets$ 密集区間を持つ単体遺伝子集合
\For{$k = 2, \ldots, k_{\max}$}
  \State $C_k \gets$ Apriori 候補生成 ($\mathcal{F}_{k-1}$)
  \State $C_k \gets$ Apriori 剪定 ($C_k$, $\mathcal{F}_{k-1}$)
  \ForAll{$S \in C_k$}
    \State 構成遺伝子の密集区間の積集合で候補範囲を限定
    \State 共起タイムスタンプを計算し、密集区間を検出
    \If{密集区間が存在}
      \State $\mathcal{F}_k \gets \mathcal{F}_k \cup \{S\}$
      \State $\mathcal{R} \gets \mathcal{R} \cup \{(S, \text{intervals})\}$
    \EndIf
  \EndFor
\EndFor
\State \Return $\mathcal{R}$
\end{algorithmic}
\end{algorithm}

\subsection{発現閾値戦略}

二値化の閾値 $\theta_j$ は結果の品質に大きく影響する。
本研究では以下の 3 戦略を比較する。

\begin{itemize}
  \item \textbf{中央値}: $\theta_j = \mathrm{median}(\{x_{ij} \mid x_{ij} > 0\})$
  \item \textbf{分位数 $q$}: $\theta_j = Q_q(\{x_{ij} \mid x_{ij} > 0\})$
  \item \textbf{z スコア}: $\theta_j = \mu_j + z \cdot \sigma_j$（非ゼロ値の統計量）
\end{itemize}

\subsection{計算量}

\begin{theorem}[計算量]
DGCI マイニングの時間計算量は $O(n \cdot |\mathcal{F}| \cdot \log n)$ である。
ここで $n$ は細胞数、$|\mathcal{F}|$ は密集区間を持つ遺伝子セット数である。
\end{theorem}

\begin{proof}
各遺伝子セットに対するスライディングウィンドウ走査は最大 $n$ 位置を訪問し、
各位置での計数は出現リスト上の二分探索により $O(\log n)$ で完了する。
Apriori の候補生成は $|\mathcal{F}|^2$ に上界されるが、
剪定により実質的な候補数は大幅に削減される。
\end{proof}

% =====================================================================
\section{実験}
\label{sec:experiments}

\subsection{実験設定}

\paragraph{合成データ}
造血分化過程を模した合成 scRNA-seq データを生成した。
$n = 500$ 細胞、25 遺伝子（造血マーカー遺伝子 15 + バックグラウンド 10）を使用し、
5 つの既知遺伝子モジュール（幹細胞維持、赤血球系コミットメント、赤血球系成熟、
骨髄系プログラム、細胞増殖）を埋め込んだ。
各モジュールは特定の擬時間区間でベル型の活性化プロファイルを持ち、
ドロップアウト率 0.3、ノイズスケール 0.5 を付加した。

\paragraph{パラメータ}
ウィンドウサイズ $W \in \{20, 30, 50, 80, 100\}$、
サポート閾値 $\sigma \in \{5, 8, 10, 15\}$、
最大アイテムセットサイズ $k_{\max} = 5$ とした。
閾値戦略は中央値、分位数（$q \in \{0.3, 0.5, 0.7\}$）、
z スコア（$z \in \{0.5, 1.0\}$）の 6 条件を比較した。

\paragraph{評価指標}
\begin{itemize}
  \item \textbf{モジュール再現率}: 少なくとも 2 遺伝子が一致する DGCI が存在する
        GT モジュールの割合
  \item \textbf{平均 IoU}: マッチした GT モジュールと検出区間の
        時間的 Intersection over Union の平均
  \item \textbf{適合率}: GT モジュールとマッチする multi-gene DGCI の割合
\end{itemize}

\subsection{E1: 感度分析}

Table~\ref{tab:sensitivity} にウィンドウサイズとサポート閾値の組み合わせによる結果を示す。

\begin{table}[t]
\centering
\caption{E1: ウィンドウサイズ $W$ とサポート閾値 $\sigma$ の感度分析（$n=500$, 分位数 $q=0.5$）}
\label{tab:sensitivity}
\begin{tabular}{rr|rrr|r}
\toprule
$W$ & $\sigma$ & 再現率 & IoU & 適合率 & Multi-gene 数 \\
\midrule
20  &  5  & 1.00 & 0.91 & 0.58 & 3724 \\
20  & 10  & 1.00 & 0.42 & 0.57 &  243 \\
50  &  8  & 1.00 & 0.92 & 0.56 & 2944 \\
50  & 15  & 1.00 & 0.88 & 0.46 &  448 \\
80  & 10  & 1.00 & 0.95 & 0.55 & 2885 \\
100 & 15  & 1.00 & 0.94 & 0.48 & 1361 \\
\bottomrule
\end{tabular}
\end{table}

全てのパラメータ設定でモジュール再現率 1.00 を達成した（$W=20, \sigma=15$ の 0.80 を除く）。
ウィンドウサイズの増大は IoU を向上させるが（$W=20$: 0.42--0.91 → $W=100$: 0.94--0.98）、
検出される DGCI 数も増加する。
サポート閾値の増大は DGCI 数を効果的に削減しつつ、高い再現率を維持する。

\subsection{E2: 閾値戦略比較}

\begin{table}[t]
\centering
\caption{E2: 発現閾値戦略の比較（$W=50$, $\sigma=8$）}
\label{tab:threshold}
\begin{tabular}{ll|rrr}
\toprule
戦略 & パラメータ & 再現率 & IoU & Multi-gene 数 \\
\midrule
中央値    & ---  & 1.00 & 0.92 & 2944 \\
分位数    & 0.3  & 1.00 & 0.94 & 12543 \\
分位数    & 0.5  & 1.00 & 0.92 & 2944 \\
分位数    & 0.7  & 1.00 & 0.87 & 790 \\
z スコア  & 0.5  & 1.00 & 0.93 & 588 \\
z スコア  & 1.0  & 1.00 & 0.96 & 169 \\
\bottomrule
\end{tabular}
\end{table}

Table~\ref{tab:threshold} に示すように、全戦略でモジュール再現率 1.00 を達成した。
分位数 $q=0.3$ は低い閾値のため最多の 12543 DGCI を検出するが、
z スコア $z=1.0$ は過発現に限定した検出を行い、最少の 169 DGCI を返す。
z スコア戦略は閾値が厳しいため検出数が少ないが、IoU は 0.96 と最も高い時間的精度を示した。

\subsection{E3: スケーラビリティ}

\begin{table}[t]
\centering
\caption{E3: 細胞数に対するスケーラビリティ}
\label{tab:scalability}
\begin{tabular}{r|rr|rr}
\toprule
$n$ (細胞数) & $W$ & $\sigma$ & Multi-gene 数 & 実行時間 (ms) \\
\midrule
 100  &  10 &   3 &  1953 &    76 \\
 200  &  20 &   4 &  3056 &   200 \\
 500  &  50 &  10 &  1498 &   296 \\
1000  & 100 &  20 &  1164 &   434 \\
2000  & 200 &  40 &  1004 &   742 \\
\bottomrule
\end{tabular}
\end{table}

Table~\ref{tab:scalability} に示すように、実行時間は細胞数に対しほぼ線形に増加し、
2000 細胞でも約 0.7 秒で完了する。
遺伝子数が固定（25 遺伝子）の条件下では、Apriori の候補空間が抑制されるため
実用的な計算時間を維持できる。

\subsection{E4: モジュール検出品質}

最適パラメータ（$W=50$, $\sigma=8$, 分位数 $q=0.4$）における
モジュール別の検出結果を Table~\ref{tab:modules} に示す。

\begin{table}[t]
\centering
\caption{E4: モジュール別検出結果}
\label{tab:modules}
\begin{tabular}{l|l|c}
\toprule
モジュール & 遺伝子セット & マッチ DGCI 数 \\
\midrule
幹細胞維持        & CD34, KIT, GATA2, RUNX1 & 709 \\
赤血球系コミットメント & GATA1, KLF1, EPOR       & 811 \\
赤血球系成熟      & HBB, HBA1, GYPA          & 513 \\
骨髄系プログラム   & CEBPA, SPI1, MPO         & 744 \\
細胞増殖          & MKI67, TOP2A, CDK1        & 665 \\
\bottomrule
\end{tabular}
\end{table}

5 つの GT モジュール全てが検出され（再現率 1.00）、
平均 IoU は 0.94 と高い時間的精度を達成した。
適合率は 0.51 であり、GT モジュールに対応しないパターンも多く検出されるが、
これは GT として定義していない遺伝子間の組み合わせも含まれるためである。

% =====================================================================
\section{考察}
\label{sec:discussion}

\paragraph{WGCNA との比較}
WGCNA は静的な共発現モジュールを検出するが、時間的局在の情報を持たない。
本手法は同じ遺伝子モジュールを検出しつつ、それが擬時間のどの区間で
活性化するかを同時に同定できる。
これは分化過程における制御プログラムの時間的順序の解明に直結する。

\paragraph{パラメータ感度}
ウィンドウサイズ $W$ とサポート閾値 $\sigma$ のバランスが重要である。
大きな $W$ は広い区間を捉えるが特異度が低下し、
大きな $\sigma$ は偽陽性を減少させるが検出力が低下する。
実データでは、擬時間の分布と生物学的文脈に基づく事前知識を活用した
パラメータ設定が推奨される。

\paragraph{制約と今後の展望}
本研究には以下の制約がある。
(1) 合成データでの評価であり、実データでの検証が必要である。
GEO や Human Cell Atlas の公開 scRNA-seq データセットへの適用が次のステップとなる。
(2) 二値化の閾値設定が結果に影響するため、データ駆動型の閾値自動選択の開発が望まれる。
(3) 分岐軌跡への拡張が必要であり、論文 C（多次元密集区間）の成果との統合が有望である。

% =====================================================================
\section{結論}
\label{sec:conclusion}

本研究では、scRNA-seq データにおける擬時間上の密集遺伝子共発現区間（DGCI）の
検出手法を提案した。
発現閾値による二値化アダプタと Apriori+スライディングウィンドウアルゴリズムにより、
従来の静的共発現解析では捉えられなかった時間的に局在した遺伝子協調パターンを同定する。
造血分化の合成データでモジュール再現率 1.00、平均 IoU 0.94 を達成し、
手法の有効性を示した。

今後は公開 scRNA-seq データセット（GEO、Human Cell Atlas）への適用、
分岐軌跡への拡張、および他の擬時間推定手法との組み合わせ評価を行う予定である。

% =====================================================================
\bibliographystyle{plain}
\bibliography{refs}

\end{CJK}
\end{document}
